\chapter{Limits and Colimits}

\begin{overviewbox}
Chapter 5 introduced the universal-property pattern through its two
simplest cases (initial and terminal objects) and previewed a richer
case (products and coproducts). This chapter develops the systematic
generalization: \emph{limits} and \emph{colimits}.

The picture is this. A \term{diagram} in a category $\mathcal{C}$ is just
a functor $D : \mathcal{J} \to \mathcal{C}$ from some small ``shape''
category $\mathcal{J}$. A \term{limit} of $D$ is a universal way to
choose a single object of $\mathcal{C}$ that ``sees the whole diagram
through the projections''; a \term{colimit} is the dual --- a universal
way to ``glue the whole diagram together along the inclusions.''

Once this template is in place, products, coproducts, equalizers,
coequalizers, pullbacks and pushouts all become instances of the same
construction with different choices of $\mathcal{J}$. Every theorem we
prove about limits will instantly give a theorem about each of these
specific cases.

We end with the slogan: \emph{limits in a functor category are computed
pointwise}. This is the first taste of why functor categories are so
well-behaved.
\end{overviewbox}


% ═══════════════════════════════════════════════════════════════════════════
\section{Diagrams as Functors}
\label{sec:diagrams-as-functors}
% ═══════════════════════════════════════════════════════════════════════════

\subsection{Informally: Picture-Then-Functor}

When we draw a diagram in $\mathcal{C}$, we are picking some objects and
some morphisms in $\mathcal{C}$ and arranging them. We do not draw every
object and every morphism; we draw a finite (or small) selection that
``has a shape.'' The shape itself is a category --- a category of
\emph{slots} or \emph{positions} for the objects, with arrows between
slots in the appropriate places.

A functor $D : \mathcal{J} \to \mathcal{C}$ then fills the slots: each
object of $\mathcal{J}$ becomes an object of $\mathcal{C}$, each arrow
becomes a morphism of $\mathcal{C}$, and the functoriality conditions
guarantee that the picture in $\mathcal{C}$ has the right relationships
between its parts.

\subsection{Expanding: The Shape Categories}

\begin{example}{Two-Object Diagram for a Product}
The shape is the discrete category on two objects, $\bullet \;\; \bullet$.
A functor from this category to $\mathcal{C}$ is a choice of two objects
of $\mathcal{C}$ --- nothing more. This is the shape we use for binary
products.
\end{example}

\begin{example}{Parallel Pair for an Equalizer}
The shape is a category with two objects and two parallel non-identity
arrows:
\[
\begin{tikzcd}
  \bullet \arrow[r, shift left] \arrow[r, shift right] & \bullet
\end{tikzcd}
\]
A functor from this category to $\mathcal{C}$ picks two objects and two
parallel morphisms $f, g : A \to B$ between them. This is the shape we
use for equalizers.
\end{example}

\begin{example}{Cospan for a Pullback}
The shape is a category with three objects arranged so that two arrows
meet at a common target:
\[
\begin{tikzcd}
  \bullet \arrow[r] & \bullet & \bullet \arrow[l]
\end{tikzcd}
\]
A functor from this category to $\mathcal{C}$ picks a cospan
$A \to C \leftarrow B$ in $\mathcal{C}$. This is the shape we use for
pullbacks.
\end{example}

\begin{dialog}
The standard moves are: pick a shape $\mathcal{J}$, then study \emph{all}
$\mathcal{J}$-shaped diagrams in $\mathcal{C}$ at once. The functor
formalism gives us a clean language for ``a $\mathcal{J}$-shaped diagram
in $\mathcal{C}$'': it is exactly an object of $[\mathcal{J},
\mathcal{C}]$.
\end{dialog}

\subsection{Formalizing: Indexing Category and Diagram}

A \term{shape} or \term{index category} is a small category $\mathcal{J}$.
A \term{diagram of shape $\mathcal{J}$ in $\mathcal{C}$} is a functor
\begin{equation*}
  D : \mathcal{J} \to \mathcal{C}.
\end{equation*}
Two such diagrams are equivalent in the sense of the functor category
$[\mathcal{J}, \mathcal{C}]$: a morphism between them is a natural
transformation.

\begin{howtosay}
\begin{itemize}
  \item ``diagram'' refers either to the picture or to the underlying
        functor $D$. We use both interchangeably.
  \item Common shapes have names: \emph{discrete} (no arrows),
        \emph{parallel pair}, \emph{cospan}, \emph{span}, \emph{walking
        arrow} (just $\bullet \to \bullet$), the natural numbers
        $\mathbb{N}$ viewed as a category, etc.
  \item The phrase ``$D$-shaped diagram'' usually means a diagram whose
        index category equals (or resembles) the shape used by $D$.
\end{itemize}
\end{howtosay}

\subsection{Simplifying: From Pictures to Functors and Back}

Every commutative-diagram picture we draw can be encoded as a functor
$D : \mathcal{J} \to \mathcal{C}$ for some $\mathcal{J}$. Going from a
picture to a functor is mechanical: the slots and arrows of the picture
\emph{are} $\mathcal{J}$, and the labels of the slots and arrows give the
action of $D$. Going from a functor to a picture is also mechanical: draw
$\mathcal{J}$ on a chalkboard, then write $D$'s values inside each slot
and on each arrow.


% ═══════════════════════════════════════════════════════════════════════════
\section{Cones and Limits}
\label{sec:cones-and-limits}
% ═══════════════════════════════════════════════════════════════════════════

\subsection{Informally: An Apex Watching the Diagram}

A \term{cone over $D$} is a single object $X$ of $\mathcal{C}$, together
with a family of morphisms --- one $\lambda_j : X \to D(j)$ for each object
$j$ of $\mathcal{J}$ --- that respects all the arrows of $\mathcal{J}$.
The word ``cone'' comes from the picture: $X$ sits as an apex above the
diagram, with rays $\lambda_j$ reaching down to each slot.

A \term{limit} of $D$ is a \emph{universal} cone: a cone $L$ such that
every other cone factors through $L$ in a unique way.

\subsection{Expanding: The Naturality Condition for Cones}

\begin{example}{Cone over a Two-Object Diagram (Product)}
A diagram of shape ``two objects, no arrows'' picks $A$ and $B$ in
$\mathcal{C}$. A cone over it is an object $X$ together with two
morphisms $X \to A$ and $X \to B$. There is no naturality constraint
because there are no non-identity arrows in the shape. So a cone is just
a pair $(f, g)$ with $f : X \to A$, $g : X \to B$. The limit, when it
exists, is the product $A \times B$ together with its projections.
\end{example}

\begin{example}{Cone over a Parallel Pair (Equalizer)}
A diagram of shape ``parallel pair'' picks $f, g : A \to B$. A cone over
it is an object $X$ with morphisms $\lambda_A : X \to A$ and $\lambda_B :
X \to B$ such that $f \circ \lambda_A = \lambda_B$ and $g \circ \lambda_A
= \lambda_B$. Combining: $f \circ \lambda_A = g \circ \lambda_A$. So
$\lambda_A$ is a morphism into $A$ that ``equalizes'' $f$ and $g$, and
$\lambda_B$ is then forced.
\end{example}

\begin{example}{Cone over a Cospan (Pullback)}
A diagram of shape $\bullet \to \bullet \leftarrow \bullet$ picks
$f : A \to C$ and $g : B \to C$. A cone over it is an object $X$ with
morphisms $\lambda_A : X \to A$, $\lambda_B : X \to B$, and
$\lambda_C : X \to C$ such that $f \circ \lambda_A = \lambda_C$ and
$g \circ \lambda_B = \lambda_C$. The morphism $\lambda_C$ is redundant
(forced by either of the others), so a pullback cone is really a pair
$(\lambda_A, \lambda_B)$ with $f \circ \lambda_A = g \circ \lambda_B$.
\end{example}

\begin{dialog}
A cone over $D$ packages a candidate limiting object $X$ with all the
``projection-like'' morphisms going down to each $D(j)$. The naturality
constraint is what makes the rays compatible with the structure of the
diagram.
\end{dialog}

\subsection{Formalizing: Cones and the Limit}

Let $\Delta_X : \mathcal{J} \to \mathcal{C}$ denote the \term{constant
functor} at $X$: every object goes to $X$, every morphism goes to
$\mathrm{id}_X$.

A \term{cone over $D$ with apex $X$} is a natural transformation
\begin{equation*}
  \lambda : \Delta_X \Rightarrow D.
\end{equation*}
Unpacking: a family $\lambda_j : X \to D(j)$ for each object $j$, such
that for every $\alpha : j \to k$ in $\mathcal{J}$,
\begin{equation*}
  D(\alpha) \circ \lambda_j = \lambda_k.
\end{equation*}
This is the naturality square for $\lambda$: it commutes because the
$\Delta_X$ side is just $\mathrm{id}_X$.

A \term{limit of $D$} is a cone $\lambda : \Delta_L \Rightarrow D$ with
the universal property: for every cone $\mu : \Delta_X \Rightarrow D$,
there is a \emph{unique} morphism $u : X \to L$ in $\mathcal{C}$ such
that $\mu_j = \lambda_j \circ u$ for every $j$.

We write $\lim D$ for the limit object (when one exists), and call the
$\lambda_j$ the \term{limit projections} or \term{limit legs}.

\begin{howtosay}
\begin{itemize}
  \item $\lim D$ --- ``limit of $D$,'' sometimes written
        $\lim_{\mathcal{J}} D$ to emphasize the indexing category.
  \item $\Delta_X$ --- ``Delta sub $X$'' or ``the constant functor at $X$.''
  \item $\lambda_j : L \to D(j)$ --- the \emph{$j$-th leg} or
        \emph{projection} of the limit cone.
  \item Limits are sometimes written $\lim_\leftarrow$ or with an
        underset arrow: the convention $\varprojlim$ also appears.
\end{itemize}
\end{howtosay}

\subsection{Reorganizing: Products, Equalizers, Pullbacks Restated}

\begin{align*}
  \text{Product of $A, B$:} & \;\; \lim D \text{ where } D : \text{discrete on } 2 \to \mathcal{C}, \\
    & \;\; D(1) = A,\; D(2) = B. \\
  \text{Equalizer of $f, g$:} & \;\; \lim D \text{ where } D : (\bullet \rightrightarrows \bullet) \to \mathcal{C}, \\
    & \;\; D \text{ sends the two arrows to } f \text{ and } g. \\
  \text{Pullback of $f, g$ over $C$:} & \;\; \lim D \text{ where } D : (\bullet \to \bullet \leftarrow \bullet) \to \mathcal{C}, \\
    & \;\; D \text{ encodes the cospan } f, g.
\end{align*}

Each of these is the same universal-property pattern of \S5.3, with a
specific shape for $\mathcal{J}$.

\subsection{Simplifying: The Limit Schema}

\begin{equation*}
  \boxed{\;
  \forall (X, \mu),\;
  \exists !\, u : X \to L,\;
  \mu_j = \lambda_j \circ u \;\;\text{for every } j.
  \;}
\end{equation*}

\subsection{Formally: Uniqueness and Functoriality of Limits}

\formalnote{Uniqueness up to unique isomorphism}
By the general theorem of \S5.4, any two limits of the same diagram are
related by a unique isomorphism compatible with the legs. We may speak
of \emph{the} limit of $D$.

\formalnote{The limit functor}
When $\mathcal{C}$ has all $\mathcal{J}$-shaped limits, the assignment
$D \mapsto \lim D$ extends to a functor
\begin{equation*}
  \lim : [\mathcal{J}, \mathcal{C}] \to \mathcal{C}.
\end{equation*}
Its action on a morphism of diagrams (a natural transformation
$\eta : D \Rightarrow D'$) sends $\lim D$ to $\lim D'$ via the unique
mediating morphism guaranteed by the limit universal property.


% ═══════════════════════════════════════════════════════════════════════════
\section{Equalizers, Pullbacks, and the Zoo of Limits}
\label{sec:limit-zoo}
% ═══════════════════════════════════════════════════════════════════════════

\subsection{Informally: One Pattern, Many Names}

We have already seen products, equalizers, and pullbacks emerge as
specific limits. This section catalogues the standard limit shapes
explicitly, so the names become familiar.

\subsection{Expanding: Each Shape, Its Universal Property}

\begin{example}{Equalizer of $f, g : A \to B$}
The equalizer is an object $E$ together with a morphism $e : E \to A$
satisfying $f \circ e = g \circ e$, universal among such pairs:
\[
\begin{tikzcd}
  E \arrow[r, "e"] & A \arrow[r, shift left, "f"] \arrow[r, shift right, "g"'] & B
\end{tikzcd}
\]
For every $h : X \to A$ with $f \circ h = g \circ h$, there is a unique
$u : X \to E$ with $e \circ u = h$.
\end{example}

\begin{example}{Pullback of $f, g$ over $C$}
The pullback (sometimes called \emph{fibered product}) is an object $P$
with morphisms $p_A : P \to A$ and $p_B : P \to B$ satisfying
$f \circ p_A = g \circ p_B$, universal among such squares:
\[
\begin{tikzcd}
  P \arrow[r, "p_B"] \arrow[d, "p_A"'] & B \arrow[d, "g"] \\
  A \arrow[r, "f"'] & C
\end{tikzcd}
\]
For every commuting square $(h, k)$ with $f \circ h = g \circ k$, there
is a unique $u : X \to P$ with $p_A \circ u = h$ and $p_B \circ u = k$.
\end{example}

\begin{example}{The Limit of $\mathbf{1}$-shaped Diagrams}
The shape $\mathbf{1}$ (one object, no non-identity morphisms) sends to
just one object $A$. A cone is a morphism $X \to A$. The limit is $A$
itself with $\lambda = \mathrm{id}_A$.
\end{example}

\begin{example}{The Limit of $\mathbf{0}$-shaped Diagrams (Terminal Object)}
The shape $\mathbf{0}$ (empty category) has no objects and no morphisms.
The functor $D : \mathbf{0} \to \mathcal{C}$ picks nothing. A cone over
$D$ is an object $X$ with no constraints --- just an object. The limit is
the \emph{terminal object} $\mathbf{1}$: the universal object with a
unique map from every other object. So terminal = limit over the empty
diagram.
\end{example}

\subsection{Formalizing: A Single Universal Property, Five Names}

For a diagram $D$ of shape $\mathcal{J}$:

\begin{align*}
  \mathcal{J} \text{ empty:} & \;\; \lim D = \mathbf{1} \;\text{(terminal object)} \\
  \mathcal{J} = \mathbf{1}:  & \;\; \lim D = D(\text{the only object}) \\
  \mathcal{J} = \text{discrete on } n: & \;\; \lim D = \prod_{i=1}^n D(i) \;\text{($n$-ary product)} \\
  \mathcal{J} = \bullet \rightrightarrows \bullet: & \;\; \lim D = \mathrm{Eq}(f, g) \;\text{(equalizer)} \\
  \mathcal{J} = \bullet \to \bullet \leftarrow \bullet: & \;\; \lim D = A \times_C B \;\text{(pullback)}
\end{align*}

\subsection{Reorganizing: The Limit Construction in $\mathbf{Set}$}

When $\mathcal{C} = \mathbf{Set}$, the limit of $D : \mathcal{J} \to
\mathbf{Set}$ has a concrete description. Define
\begin{equation*}
  \lim D = \left\{ (x_j)_{j \in \mathrm{Ob}(\mathcal{J})}
  \;\Big|\; x_j \in D(j),\;
  D(\alpha)(x_j) = x_k \text{ for every } \alpha : j \to k \right\}.
\end{equation*}
The legs $\lambda_j$ are the obvious projections. One can check this
satisfies the universal property. So in $\mathbf{Set}$ ``limit'' is
``compatible tuples.''

\subsection{Simplifying: When Do Limits Exist?}

A category has \emph{finite limits} if it has limits for every diagram
whose shape $\mathcal{J}$ has finitely many objects and morphisms.
Equivalently (the limit zoo collapses): a category has finite limits iff
it has a terminal object, all binary products, and all equalizers.

A category has \emph{all (small) limits} if it has limits for every
small-shaped diagram. Such a category is called \term{(small-)complete}.

\subsection{Formally: The Equalizer Construction in $\mathbf{Set}$}

\formalnote{Equalizers in $\mathbf{Set}$}
For $f, g : A \to B$ in $\mathbf{Set}$:
\begin{equation*}
  \mathrm{Eq}(f, g) = \{a \in A \mid f(a) = g(a)\},\;\;
  e : \mathrm{Eq}(f, g) \hookrightarrow A.
\end{equation*}
Pullbacks have a similar concrete description:
\begin{equation*}
  A \times_C B = \{(a, b) \in A \times B \mid f(a) = g(b)\}.
\end{equation*}


% ═══════════════════════════════════════════════════════════════════════════
\section{Colimits as Dual Limits}
\label{sec:colimits}
% ═══════════════════════════════════════════════════════════════════════════

\subsection{Informally: Glue Instead of Project}

The dual of a limit is a colimit. Where a limit is an apex \emph{above}
the diagram with rays going \emph{down}, a colimit is a vertex \emph{below}
the diagram with rays going \emph{up}. Where a limit's universal property
is about \emph{mapping in}, a colimit's is about \emph{mapping out}.

Roughly, a colimit \emph{glues together} the objects in the diagram along
the morphisms in the diagram. In $\mathbf{Set}$, this gluing is
realized as a quotient of a disjoint union by an equivalence relation.

\subsection{Expanding: Standard Colimit Shapes}

\begin{example}{Coproducts}
For the discrete two-object shape: the colimit is the coproduct $A + B$
together with inclusions $\iota_1, \iota_2$.
\end{example}

\begin{example}{Coequalizers}
For the parallel-pair shape with $f, g : A \to B$: the colimit is the
\term{coequalizer} of $f$ and $g$ --- the universal object $Q$ with a
morphism $q : B \to Q$ such that $q \circ f = q \circ g$. In $\mathbf{Set}$,
$Q$ is $B$ modulo the equivalence relation generated by $f(a) \sim g(a)$
for all $a \in A$.
\end{example}

\begin{example}{Pushouts}
For the span shape $\bullet \leftarrow \bullet \to \bullet$ with
$f : C \to A$ and $g : C \to B$: the colimit is the \term{pushout},
written $A +_C B$ or $A \sqcup_C B$. In $\mathbf{Set}$, it is
$(A \sqcup B) / \sim$, where $f(c) \sim g(c)$ for all $c \in C$.
\end{example}

\subsection{Formalizing: Cocones and Colimits}

A \term{cocone under $D$ with co-apex $X$} is a natural transformation
\begin{equation*}
  \mu : D \Rightarrow \Delta_X.
\end{equation*}
Unpacking: morphisms $\mu_j : D(j) \to X$ for each $j$, with
$\mu_k \circ D(\alpha) = \mu_j$ for every $\alpha : j \to k$.

A \term{colimit of $D$} is a universal cocone: $\mu : D \Rightarrow
\Delta_C$ such that every cocone $\nu : D \Rightarrow \Delta_X$ factors
uniquely as $\nu_j = u \circ \mu_j$ for some unique $u : C \to X$.

We write $\mathrm{colim}\, D$ (also $\lim_\to$ or $\varinjlim D$) for the
colimit.

\begin{howtosay}
\begin{itemize}
  \item $\mathrm{colim}\, D$ --- ``colimit of $D$''
  \item $\varinjlim D$ or $\lim_\to$ --- the arrow-direction notation
  \item ``cocone'' is sometimes written ``co-cone''
  \item Specific names: coproduct, coequalizer, pushout, initial object
        (colimit over empty)
\end{itemize}
\end{howtosay}

\subsection{Reorganizing: Duality, Once Stated}

A colimit of $D : \mathcal{J} \to \mathcal{C}$ is a limit of
$D^{\mathrm{op}} : \mathcal{J}^{\mathrm{op}} \to \mathcal{C}^{\mathrm{op}}$.
Therefore every theorem we prove about limits has a dual for colimits, by
reversing arrows. We will prove things about limits and rely on duality
for the colimit side without comment.

\subsection{Simplifying: The Colimit Schema}

\begin{equation*}
  \boxed{\;
  \forall (X, \nu),\;
  \exists !\, u : C \to X,\;
  \nu_j = u \circ \mu_j \;\;\text{for every } j.
  \;}
\end{equation*}

\subsection{Formally: Cocomplete Categories}

\formalnote{Cocompleteness}
A category is \term{cocomplete} if it has all small colimits. Equivalently,
it has an initial object, all small coproducts, and all coequalizers.

\formalnote{The free constructions}
Many ``free'' constructions in mathematics --- free monoid on a set, free
group on a set, free $R$-module on a set --- are colimits in disguise.
Free constructions go hand-in-hand with adjunctions (Chapter~7).


% ═══════════════════════════════════════════════════════════════════════════
\section{Preservation and Reflection of Limits}
\label{sec:preservation}
% ═══════════════════════════════════════════════════════════════════════════

\subsection{Informally: Which Functors Behave Nicely?}

Once we have limits and colimits, the natural question is: which functors
\emph{respect} them? A functor that respects limits is invaluable, because
we can compute the limit in the source and then apply the functor, rather
than computing inside the target.

Three properties matter: \emph{preservation} (the functor sends limits to
limits), \emph{reflection} (the functor turns limits-in-target back into
limits-in-source), and \emph{creation} (the strongest, which we discuss
formally).

\subsection{Expanding: A Functor That Preserves Products}

\begin{example}{The Forgetful Functor $U : \mathbf{Mon} \to \mathbf{Set}$}
$U$ sends a monoid to its underlying set, and a monoid homomorphism to
the underlying function. It preserves products: $U(M_1 \times M_2) =
U(M_1) \times U(M_2)$, with projections sent to projections. We will see
in Chapter~7 that $U$ is the right adjoint of a free-monoid functor, and
right adjoints preserve all limits.
\end{example}

\begin{example}{A Functor That Does Not Preserve Limits}
The functor $\mathbf{Set} \to \mathbf{Set}$ that sends $X$ to the
two-element set $\{0, 1\}$ regardless of $X$ does not preserve products:
$\{0,1\} = F(A \times B) \neq F(A) \times F(B) = \{0,1\} \times \{0,1\}$
which has four elements. Constant functors typically do not preserve
limits.
\end{example}

\subsection{Formalizing: Preservation, Reflection, Creation}

Let $F : \mathcal{C} \to \mathcal{D}$ be a functor, $D : \mathcal{J} \to
\mathcal{C}$ a diagram, and $\lambda : \Delta_L \Rightarrow D$ a limiting
cone over $D$.

\begin{itemize}
  \item $F$ \term{preserves} the limit of $D$ if $F \lambda :
        \Delta_{F(L)} \Rightarrow F \circ D$ is a limiting cone over
        $F \circ D$.
  \item $F$ \term{reflects} the limit of $D$ if: whenever
        $\lambda' : \Delta_X \Rightarrow D$ is a cone over $D$ in
        $\mathcal{C}$ and $F \lambda'$ is limiting over $F \circ D$,
        then $\lambda'$ was limiting over $D$ already.
  \item $F$ \term{creates} the limit of $D$ if every limiting cone over
        $F \circ D$ in $\mathcal{D}$ uniquely lifts to a cone over $D$
        in $\mathcal{C}$, and the lifted cone is itself a limit of $D$.
\end{itemize}

Preservation says ``limits in $\mathcal{C}$ map to limits in $\mathcal{D}$.''
Reflection says ``limits in $\mathcal{D}$ come from limits in $\mathcal{C}$.''
Creation says ``the limit in $\mathcal{C}$ is determined by the limit in
$\mathcal{D}$.''

\subsection{Reorganizing: A Slogan and Two Anchors}

\textbf{Anchor 1.}\quad Right adjoints preserve limits (Chapter~7).

\textbf{Anchor 2.}\quad Fully faithful functors reflect limits.

These two facts together account for an enormous fraction of the
limit-preservation phenomena one meets in practice.

\subsection{Simplifying: The Preservation Schema}

\begin{equation*}
  F \text{ preserves } \lim D
  \;\iff\;
  F(\lim D) \cong \lim (F \circ D),
\end{equation*}
\emph{naturally and via the canonical mediating morphism}, not merely up
to abstract isomorphism.

\subsection{Formally: Limits in Functor Categories}

\formalnote{Limits are computed pointwise}
For a small category $\mathcal{J}$ and any category $\mathcal{D}$ with
$\mathcal{J}$-limits, the functor category $[\mathcal{C}, \mathcal{D}]$
has $\mathcal{J}$-limits, and they are computed pointwise:
\begin{equation*}
  (\lim D)(C) \cong \lim_{j} D(j)(C)
\end{equation*}
for each object $C$ of $\mathcal{C}$. The evaluation functors
$\mathrm{ev}_C : [\mathcal{C}, \mathcal{D}] \to \mathcal{D}$ jointly create
all limits. This is one of the most-used facts in category theory:
calculations in functor categories reduce to calculations object-by-object
in the target.


% ═══════════════════════════════════════════════════════════════════════════
\section*{Exercises}
\addcontentsline{toc}{section}{Exercises}
% ═══════════════════════════════════════════════════════════════════════════

\begin{exercisesbox}
\begin{qa}
\qa{What is a diagram in $\mathcal{C}$, formally?}{A functor $D : \mathcal{J} \to \mathcal{C}$ from a small category $\mathcal{J}$.}
\qa{What plays the role of ``the shape''?}{The index category $\mathcal{J}$.}
\qa{What is the constant functor $\Delta_X$?}{The functor sending every object to $X$ and every arrow to $\mathrm{id}_X$.}
\qa{What is a cone over $D$ with apex $X$?}{A natural transformation $\Delta_X \Rightarrow D$.}
\qa{Unpacked, what data does a cone consist of?}{A morphism $\lambda_j : X \to D(j)$ for every $j$, compatible with the arrows of $\mathcal{J}$.}
\qa{What is a limit of $D$?}{A universal cone over $D$.}
\qa{Why is the limit unique up to unique iso?}{It satisfies a universal property; \S5.4 applies.}
\qa{Empty index category: what is the limit?}{The terminal object.}
\qa{Singleton index category: what is the limit?}{The unique object that $D$ picks out.}
\qa{Discrete two-object index: what is the limit?}{The (binary) product.}
\qa{Parallel pair index: what is the limit?}{The equalizer.}
\qa{Cospan index: what is the limit?}{The pullback.}
\qa{What is a cocone under $D$?}{A natural transformation $D \Rightarrow \Delta_X$.}
\qa{What is a colimit?}{A universal cocone.}
\qa{Empty index: what is the colimit?}{The initial object.}
\qa{Discrete two-object index: colimit?}{The coproduct.}
\qa{Parallel pair: colimit?}{The coequalizer.}
\qa{Span index: colimit?}{The pushout.}
\qa{Why call the colimit ``a glue''?}{It glues the diagram's objects together along the diagram's arrows.}
\qa{In $\mathbf{Set}$, how is the pullback realized?}{As $\{(a,b) \in A \times B \mid f(a) = g(b)\}$.}
\qa{In $\mathbf{Set}$, how is the equalizer realized?}{As $\{a \in A \mid f(a) = g(a)\}$ with the inclusion into $A$.}
\qa{In $\mathbf{Set}$, how is the coequalizer realized?}{As $B$ modulo the equivalence relation generated by $f(a) \sim g(a)$.}
\qa{What does ``$\mathcal{C}$ is complete'' mean?}{$\mathcal{C}$ has all small limits.}
\qa{What is the minimal checklist for completeness?}{Terminal object, all binary products, all equalizers.}
\qa{What is ``$\mathcal{C}$ is cocomplete''?}{Has all small colimits.}
\qa{Minimal checklist for cocompleteness?}{Initial object, all binary coproducts, all coequalizers.}
\qa{What does ``$F$ preserves limits'' mean?}{$F$ sends limiting cones to limiting cones.}
\qa{What is ``$F$ reflects limits''?}{If $F\lambda$ is a limit, then $\lambda$ was already a limit.}
\qa{Where are limits computed in $[\mathcal{C}, \mathcal{D}]$?}{Pointwise, in $\mathcal{D}$.}
\qa{What is the formula for $(\lim D)(C)$ pointwise?}{$\lim_{j} D(j)(C)$.}
\qa{Why is ``pointwise'' so useful?}{It reduces hard calculations in functor categories to easier ones in the target category.}
\qa{Are products limits over what shape?}{The discrete category on two (or more) objects.}
\qa{Why is duality a labour-saver here?}{We prove a theorem about limits, and get a colimit theorem free by reversing arrows.}
\end{qa}
\end{exercisesbox}


% ═══════════════════════════════════════════════════════════════════════════
\section*{Chapter Summary}
\addcontentsline{toc}{section}{Chapter Summary}
% ═══════════════════════════════════════════════════════════════════════════

\begin{summarybox}
A \textbf{diagram} in $\mathcal{C}$ is a functor $D : \mathcal{J} \to
\mathcal{C}$ from a small shape category. A \textbf{cone} over $D$ is a
natural transformation $\Delta_X \Rightarrow D$; a \textbf{cocone} under
$D$ is $D \Rightarrow \Delta_X$.

A \textbf{limit} of $D$ is a universal cone; a \textbf{colimit} is a
universal cocone. Products, equalizers, pullbacks (and their duals)
are all limits/colimits of specific shapes.

Limits and colimits are unique up to unique isomorphism, and they are
\textbf{computed pointwise} in functor categories.

A category is \textbf{complete} if it has all small limits, and
\textbf{cocomplete} if it has all small colimits. Equivalently:
terminal/initial + binary products/coproducts + equalizers/coequalizers.

A functor \textbf{preserves} a limit when it sends limiting cones to
limiting cones; it \textbf{reflects} when target-limits descend to
source-limits.
\end{summarybox}


% ═══════════════════════════════════════════════════════════════════════════
\section*{Formal Restatement}
\addcontentsline{toc}{section}{Formal Restatement}
% ═══════════════════════════════════════════════════════════════════════════

\begin{formalrestatement}
\textbf{Diagram.}\quad
A diagram of shape $\mathcal{J}$ in $\mathcal{C}$ is a functor
$D : \mathcal{J} \to \mathcal{C}$.

\medskip
\textbf{Cone.}\quad
A cone over $D$ with apex $X$ is a natural transformation
$\lambda : \Delta_X \Rightarrow D$:
\begin{equation*}
  \lambda_j : X \to D(j),\;\;
  D(\alpha) \circ \lambda_j = \lambda_k
  \;\text{for every } \alpha : j \to k.
\end{equation*}

\medskip
\textbf{Limit.}\quad
$L$ is a limit of $D$, with cone $\lambda$, iff for every cone
$\mu : \Delta_X \Rightarrow D$:
\begin{equation*}
  \exists !\, u : X \to L \;\text{with}\;
  \mu_j = \lambda_j \circ u \text{ for all } j.
\end{equation*}

\medskip
\textbf{Colimit.}\quad
$C$ is a colimit of $D$, with cocone $\mu$, iff for every cocone
$\nu : D \Rightarrow \Delta_X$:
\begin{equation*}
  \exists !\, u : C \to X \;\text{with}\;
  \nu_j = u \circ \mu_j \text{ for all } j.
\end{equation*}

\medskip
\textbf{Completeness.}\quad
$\mathcal{C}$ is (small-)complete iff it has all small limits, iff it has
a terminal object, all binary products, and all equalizers. Dually for
cocompleteness.

\medskip
\textbf{Pointwise limits.}\quad
$[\mathcal{C}, \mathcal{D}]$ has $\mathcal{J}$-limits whenever
$\mathcal{D}$ does, and:
\begin{equation*}
  (\lim_j D_j)(C) = \lim_j (D_j(C)).
\end{equation*}

\medskip
\noindent\textit{Chapter~7 develops adjunctions, the most important
source of limit-preserving and colimit-preserving functors.}
\end{formalrestatement}
